\documentclass[11pt,aspectratio=169]{beamer}
\usetheme{SimplePlus}

\usepackage[T1]{fontenc}
\usepackage[utf8]{inputenc}
\usepackage{amsmath,amssymb,mathtools}
\usepackage{booktabs}
\usepackage{array}
\usepackage{appendixnumberbeamer}
\usepackage{hyperref}

\newcommand{\E}{\mathbb{E}}
\newcommand{\Prob}{\mathbb{P}}
\newcommand{\Var}{\mathrm{Var}}
\newcommand{\Cov}{\mathrm{Cov}}
\newcommand{\plim}{\mathop{\mathrm{plim}}}
\newcommand{\ATE}{\mathrm{ATE}}
\newcommand{\LATE}{\mathrm{LATE}}
\newcommand{\ITT}{\mathrm{ITT}}
\newcommand{\indep}{\perp\!\!\!\perp}
\newcommand{\derivbutton}[1]{\vspace{0.4em}\hfill\hyperlink{app:#1}{\beamerbutton{Appendix details}}}
\newcommand{\backbutton}[1]{\vspace{0.4em}\hfill\hyperlink{main:#1}{\beamerbutton{Back to main slide}}}

\title{Applied Econometrics: Session 5}
\subtitle{Instrumental Variables, Non-Compliance, and LATE}
\author{PhD Intensive Course}
\institute{Lecture 5}
\date{\today}

\begin{document}

\begin{frame}
  \titlepage
\end{frame}

\begin{frame}{Why This Lecture Matters}
  Lectures 1--4 taught one recurring lesson: comparisons are causal only when the missing counterfactual is credible.
  \vspace{0.4em}

  Today we study a different route to credibility:
  \begin{block}{Core idea}
    Use a variable \(Z_i\) that shifts treatment \(D_i\), but is otherwise unrelated to potential outcomes.
  \end{block}

  \begin{itemize}
    \item \(D_i\) may be selected on unobserved ability, preferences, income, or health.
    \item \(Z_i\) supplies cleaner variation in \(D_i\).
    \item The price is that IV estimates a special effect: often the effect for compliers.
  \end{itemize}
\end{frame}

\begin{frame}{Roadmap for Today}
  \tableofcontents
\end{frame}

\section{Why IV?}

\begin{frame}{The Problem: A Hidden Confounder}
  Suppose we want the effect of schooling on wages:
  \[
    Y_i = \alpha + \beta D_i + u_i,
  \]
  where \(Y_i\) is log wage and \(D_i\) is years of schooling.

  \begin{itemize}
    \item If ability is omitted, it is inside \(u_i\).
    \item Ability likely affects both schooling and wages.
    \item Then \(\Cov(D_i,u_i)\neq 0\), so OLS does not isolate \(\beta\).
  \end{itemize}

  \begin{block}{Translation}
    We are comparing people with different education, but also different unobserved traits.
  \end{block}
\end{frame}

\begin{frame}{Why More Controls May Not Save Us}
  The ideal long regression is
  \[
    Y_i = \alpha + \beta D_i + W_i'\gamma + \varepsilon_i,
    \qquad
    \E[\varepsilon_i\mid D_i,W_i]=0.
  \]
  But sometimes \(W_i\) is missing or badly measured.

  \begin{itemize}
    \item Ability, motivation, parental inputs, and networks are hard to observe.
    \item Proxy controls can reduce bias, but they rarely remove it mechanically.
    \item IV asks for a different object: a source of variation in \(D_i\) that is not chosen by the individual.
  \end{itemize}

  \begin{block}{Practical lesson}
    IV is not a better regression control. It is a design argument about where treatment variation comes from.
  \end{block}
\end{frame}

\begin{frame}{The IV Picture}
  Think of three variables:
  \[
    Z_i \longrightarrow D_i \longrightarrow Y_i.
  \]
  A valid instrument has two roles:
  \begin{enumerate}
    \item It moves treatment: \(Z_i\) changes \(D_i\).
    \item It is clean: \(Z_i\) has no other route to \(Y_i\).
  \end{enumerate}

  \vspace{0.3em}
  A hidden confounder \(U_i\) may still affect both \(D_i\) and \(Y_i\):
  \[
    U_i \longrightarrow D_i,
    \qquad
    U_i \longrightarrow Y_i.
  \]

  \begin{block}{IV logic}
    We do not need \(D_i\) to be as good as random. We need the part of \(D_i\) caused by \(Z_i\) to be as good as random.
  \end{block}
\end{frame}

\begin{frame}{Examples of Instruments}
  \begin{itemize}
    \item \textbf{Quarter of birth and schooling:} compulsory schooling laws make school-leaving age depend slightly on birth timing.
    \item \textbf{Draft lottery and military service:} random lottery number changes the probability of service.
    \item \textbf{Random offer and program take-up:} an encouragement changes who joins a program.
    \item \textbf{Distance to college and education:} proximity can reduce cost of schooling, but exclusion is controversial.
  \end{itemize}

  \begin{block}{Student warning}
    ``A variable that predicts treatment'' is not enough. Many predictors of treatment are also predictors of potential outcomes.
  \end{block}
\end{frame}

\begin{frame}{Instrument or Control?}
  A control variable is something we condition on.
  \[
    Y_i = \alpha + \beta D_i + X_i'\gamma + e_i.
  \]
  An instrument is something we use to generate treatment variation.
  \[
    D_i = \pi_0 + \pi_1 Z_i + X_i'\pi_2 + v_i.
  \]

  \begin{itemize}
    \item Controls block back-door paths from \(D_i\) to \(Y_i\).
    \item Instruments create front-door-like variation from \(Z_i\) to \(D_i\).
    \item You can use controls inside an IV design, but the identifying variation still comes from \(Z_i\).
  \end{itemize}
\end{frame}

\section{Valid Instruments}

\begin{frame}{Notation for the Basic IV Model}
  \hypertarget{main:iv-conditions}{}
  Start with
  \[
    Y_i = \alpha + \beta D_i + u_i,
    \qquad
    \Cov(D_i,u_i)\neq 0.
  \]
  We observe an instrument \(Z_i\).

  \begin{block}{Three core conditions}
    \begin{enumerate}
      \item \textbf{Relevance:} \(\Cov(Z_i,D_i)\neq 0\).
      \item \textbf{Independence:} \(Z_i\) is as good as randomly assigned relative to potential outcomes and unobservables.
      \item \textbf{Exclusion:} \(Z_i\) affects \(Y_i\) only through \(D_i\).
    \end{enumerate}
  \end{block}

  \derivbutton{iv-moment}
\end{frame}

\begin{frame}{Condition 1: Relevance}
  Relevance means the instrument changes treatment:
  \[
    \Cov(Z_i,D_i)\neq 0.
  \]

  Estimate the first stage:
  \[
    D_i = \pi_0 + \pi_1 Z_i + r_i.
  \]

  \begin{itemize}
    \item If \(\pi_1=0\), the instrument gives no information about treatment.
    \item If \(\pi_1\) is tiny, IV estimates are unstable and noisy.
    \item Relevance can be checked in the data.
  \end{itemize}

  \begin{block}{Classroom phrase}
    The first stage asks: does the instrument actually move the treatment?
  \end{block}
\end{frame}

\begin{frame}{Condition 2: Independence}
  Independence means \(Z_i\) is not chosen because of the same unobserved forces that determine \(Y_i\).

  A strong version is
  \[
    Z_i \indep \{Y_i(d),D_i(z): d,z\in\{0,1\}\}.
  \]

  \begin{itemize}
    \item Random assignment of encouragement often makes this plausible.
    \item Institutional rules can make this plausible, but not automatic.
    \item Balance checks help detect problems, but cannot prove independence.
  \end{itemize}

  \begin{block}{Interpretation}
    Independence is about who gets the instrument, not who gets the treatment.
  \end{block}
\end{frame}

\begin{frame}{Condition 3: Exclusion Restriction}
  Exclusion says the instrument has no direct effect on the outcome.
  \[
    Z_i \rightarrow D_i \rightarrow Y_i,
    \qquad
    \text{but not}
    \qquad
    Z_i \rightarrow Y_i.
  \]

  In potential-outcome notation:
  \[
    Y_i(d,z)=Y_i(d)
    \quad \text{for all } d,z.
  \]

  \begin{itemize}
    \item A schooling instrument must not affect wages except through schooling.
    \item A program encouragement must not affect outcomes through information, attention, or stigma except through participation.
    \item Exclusion is usually argued from institutional knowledge, not tested directly.
  \end{itemize}
\end{frame}

\begin{frame}{What Can Be Tested?}
  \begin{center}
  \begin{tabular}{p{0.28\linewidth}p{0.26\linewidth}p{0.34\linewidth}}
    \toprule
    Requirement & Data check? & What to do in practice \\
    \midrule
    Relevance & Yes & Report first-stage coefficient and \(F\)-statistic \\
    Independence & Partly & Check balance on pre-treatment covariates \\
    Exclusion & No direct test & Defend the institutional story; discuss threats \\
    Monotonicity & Usually no direct test & Argue no one moves opposite to the instrument \\
    \bottomrule
  \end{tabular}
  \end{center}

  \begin{block}{Exam point}
    A significant first stage does not validate the instrument. It validates only one required condition.
  \end{block}
\end{frame}

\begin{frame}{Bad Instruments: Common Failure Modes}
  \begin{itemize}
    \item \textbf{Weak instrument:} \(Z_i\) barely moves \(D_i\).
    \item \textbf{Direct effect:} \(Z_i\) affects \(Y_i\) through a path other than treatment.
    \item \textbf{Selection into instrument:} people with different potential outcomes are more likely to receive \(Z_i\).
    \item \textbf{Multiple treatments:} \(Z_i\) changes several behaviors, not just the named treatment.
  \end{itemize}

  \begin{block}{Practical rule}
    Always write the story as a causal graph in words. If you cannot explain the path from \(Z_i\) to \(D_i\) and rule out other paths to \(Y_i\), do not trust IV.
  \end{block}
\end{frame}

\section{Wald IV}

\begin{frame}{The First Stage}
  \hypertarget{main:first-stage}{}
  The first stage estimates how treatment changes with the instrument:
  \[
    D_i = \pi_0 + \pi_1 Z_i + r_i.
  \]
  For binary \(Z_i\),
  \[
    \pi_1 = \E[D_i\mid Z_i=1]-\E[D_i\mid Z_i=0].
  \]

  \begin{itemize}
    \item If \(D_i\) is schooling, the first stage is extra years of schooling caused by \(Z_i\).
    \item If \(D_i\) is program take-up, the first stage is the compliance rate induced by assignment.
  \end{itemize}

  \begin{block}{Units}
    First stage has units of treatment per unit of instrument.
  \end{block}
\end{frame}

\begin{frame}{The Reduced Form}
  The reduced form estimates how the outcome changes with the instrument:
  \[
    Y_i = \rho_0 + \rho_1 Z_i + \eta_i.
  \]
  For binary \(Z_i\),
  \[
    \rho_1 = \E[Y_i\mid Z_i=1]-\E[Y_i\mid Z_i=0].
  \]

  \begin{itemize}
    \item The reduced form is an effect of the instrument on the outcome.
    \item Under exclusion, that effect must operate through treatment.
    \item It is often called an intention-to-treat effect in encouragement designs.
  \end{itemize}

  \begin{block}{Units}
    Reduced form has units of outcome per unit of instrument.
  \end{block}
\end{frame}

\begin{frame}{The Wald Estimator}
  \hypertarget{main:wald}{}
  If \(Z_i\) is binary, the IV estimator is the ratio:
  \[
    \hat{\beta}_{Wald}
    =
    \frac{\bar{Y}_{Z=1}-\bar{Y}_{Z=0}}
         {\bar{D}_{Z=1}-\bar{D}_{Z=0}}.
  \]

  Equivalently,
  \[
    \beta_{IV}
    =
    \frac{\Cov(Y_i,Z_i)}{\Cov(D_i,Z_i)}
    =
    \frac{\text{reduced form}}{\text{first stage}}.
  \]

  \begin{block}{Intuition}
    Convert the effect of the instrument on the outcome into outcome units per one unit of treatment.
  \end{block}

  \derivbutton{wald-proof}
\end{frame}

\begin{frame}{A Numerical Wald Example}
  Suppose assignment to a tutoring offer changes tutoring take-up and exam scores:
  \begin{center}
  \begin{tabular}{lcc}
    \toprule
    & Z=1 & Z=0 \\
    \midrule
    \(\E[D\mid Z]\) & 0.70 & 0.30 \\
    \(\E[Y\mid Z]\) & 74 & 70 \\
    \bottomrule
  \end{tabular}
  \end{center}

  First stage:
  \[
    0.70-0.30=0.40.
  \]
  Reduced form:
  \[
    74-70=4.
  \]
  Wald IV:
  \[
    \frac{4}{0.40}=10.
  \]

  Interpretation: among the students whose tutoring status was changed by the offer, tutoring raised scores by 10 points.
\end{frame}

\begin{frame}{Why the Ratio Makes Sense}
  Imagine the offer raises average test scores by 4 points.

  \begin{itemize}
    \item But only 40 percentage points of students change tutoring status because of the offer.
    \item The observed 4-point difference is diluted by students whose treatment did not change.
    \item Dividing by 0.40 undilutes the assignment effect into a treatment effect.
  \end{itemize}

  \begin{block}{Careful language}
    This is not the effect for everyone unless treatment effects are constant or the complier effect equals the population effect.
  \end{block}
\end{frame}

\begin{frame}{IV with Controls}
  In many applications we include controls \(X_i\):
  \[
    Y_i = \alpha + \beta D_i + X_i'\gamma + u_i.
  \]
  The instrument must be valid conditional on those controls.

  \begin{itemize}
    \item First stage: regress \(D_i\) on \(Z_i\) and \(X_i\).
    \item Reduced form: regress \(Y_i\) on \(Z_i\) and \(X_i\).
    \item IV ratio uses the residual variation in \(Z_i\) after removing \(X_i\).
  \end{itemize}

  \begin{block}{Rule}
    Any exogenous control in the outcome equation belongs in the first stage too.
  \end{block}

  \derivbutton{controls-proof}
\end{frame}

\section{2SLS}

\begin{frame}{Two-Stage Least Squares: Stage 1}
  \hypertarget{main:2sls}{}
  Stage 1 predicts treatment using the instrument and controls:
  \[
    D_i = \pi_0 + \pi_1 Z_i + X_i'\pi_2 + r_i.
  \]
  Get fitted values:
  \[
    \hat{D}_i = \hat{\pi}_0 + \hat{\pi}_1 Z_i + X_i'\hat{\pi}_2.
  \]

  \begin{itemize}
    \item \(\hat{D}_i\) is the part of treatment explained by \(Z_i\) and \(X_i\).
    \item The key identifying variation in \(\hat{D}_i\) comes from \(Z_i\).
  \end{itemize}
\end{frame}

\begin{frame}{Two-Stage Least Squares: Stage 2}
  Stage 2 regresses the outcome on predicted treatment and controls:
  \[
    Y_i = \alpha + \beta \hat{D}_i + X_i'\gamma + \text{error}_i.
  \]

  \begin{itemize}
    \item The coefficient on \(\hat{D}_i\) is the 2SLS estimate.
    \item With one endogenous variable and one instrument, it equals the Wald ratio after residualizing controls.
    \item Software must compute the standard errors; manual second-stage OLS standard errors are wrong.
  \end{itemize}

  \derivbutton{2sls-proof}
\end{frame}

\begin{frame}{What 2SLS Is Doing Geometrically}
  OLS uses all variation in \(D_i\):
  \[
    D_i = \text{clean variation} + \text{confounded variation}.
  \]
  2SLS keeps the variation in \(D_i\) that is explained by the instrument:
  \[
    \hat{D}_i = \mathrm{Projection}(D_i \mid Z_i,X_i).
  \]

  \begin{block}{Intuition}
    If \(Z_i\) is clean, then the instrument-generated part of treatment is clean.
  \end{block}

  \begin{itemize}
    \item 2SLS does not make every treated and untreated person comparable.
    \item It compares people whose treatment differs because of \(Z_i\).
  \end{itemize}
\end{frame}

\begin{frame}{Multiple Instruments}
  If there are several instruments \(Z_i=(Z_{1i},Z_{2i},\ldots,Z_{ki})\), stage 1 is
  \[
    D_i = \pi_0 + Z_i'\pi_1 + X_i'\pi_2 + r_i.
  \]

  \begin{itemize}
    \item Multiple instruments can strengthen prediction of treatment.
    \item They also multiply the number of exclusion restrictions to defend.
    \item With too many weak or questionable instruments, 2SLS can move toward the biased OLS estimate.
  \end{itemize}

  \begin{block}{Practical rule}
    More instruments are not automatically better. Each instrument needs its own validity argument.
  \end{block}
\end{frame}

\begin{frame}{Do Not Run IV by Hand for Inference}
  Manually running the two regressions is useful for intuition:
  \[
    D_i \rightarrow \hat{D}_i,
    \qquad
    Y_i \text{ on } \hat{D}_i.
  \]
  But the second-stage residuals are not ordinary OLS residuals.

  \begin{itemize}
    \item The second stage uses an estimated regressor.
    \item Naive OLS standard errors ignore first-stage uncertainty.
    \item Use an IV/2SLS routine for final tables.
  \end{itemize}

  \begin{block}{For this course}
    By hand for understanding; proper IV software for reported inference.
  \end{block}
\end{frame}

\section{Weak Instruments}

\begin{frame}{Weak Instruments}
  \hypertarget{main:weak}{}
  A weak instrument has a small first stage:
  \[
    D_i = \pi_0 + \pi_1 Z_i + r_i,
    \qquad
    \pi_1 \approx 0.
  \]

  \begin{itemize}
    \item The denominator of the IV ratio is small.
    \item Small sampling noise in the reduced form gets magnified.
    \item Confidence intervals become wide and conventional approximations can be misleading.
  \end{itemize}

  \begin{block}{Visual intuition}
    Weak IV is like estimating a slope by dividing by a nearly flat first-stage line.
  \end{block}

  \derivbutton{weak-proof}
\end{frame}

\begin{frame}{The First-Stage \(F\)-Statistic}
  For one endogenous treatment, a common diagnostic is the first-stage \(F\)-statistic on the excluded instruments.

  \begin{itemize}
    \item Very small \(F\): the instrument barely shifts treatment.
    \item Larger \(F\): more information about treatment from the instrument.
    \item A classic rule of thumb is \(F>10\), but it is not a theorem of validity.
  \end{itemize}

  \begin{block}{Important distinction}
    Strong first stage helps precision and weak-IV bias. It does not prove exclusion or independence.
  \end{block}
\end{frame}

\begin{frame}{Weak-IV Symptoms}
  In empirical work, weak instruments often produce:
  \begin{itemize}
    \item large standard errors,
    \item estimates that change a lot across specifications,
    \item implausibly large point estimates from tiny first stages,
    \item IV estimates drifting toward OLS in finite samples,
    \item confidence intervals that are not trustworthy under standard approximations.
  \end{itemize}

  \begin{block}{Practical response}
    Report the first stage clearly, show robustness, and be honest if the design has little treatment variation.
  \end{block}
\end{frame}

\begin{frame}{What IV Does Not Fix}
  IV does not solve every causal problem.

  \begin{itemize}
    \item It does not fix a bad exclusion restriction.
    \item It does not identify effects for people whose treatment is unaffected by the instrument.
    \item It does not make finite samples behave like large samples.
    \item It does not remove the need for a clear research question.
  \end{itemize}

  \begin{block}{Bottom line}
    IV trades one impossible requirement, observing all confounders, for another demanding requirement, finding valid instrument variation.
  \end{block}
\end{frame}

\section{Non-Compliance}

\begin{frame}{Assignment Versus Treatment Received}
  \hypertarget{main:noncompliance}{}
  Many experiments randomize an assignment \(Z_i\), but people do not perfectly follow it.

  \[
    Z_i = \text{assigned treatment},
    \qquad
    D_i = \text{received treatment}.
  \]

  \begin{itemize}
    \item Some assigned to treatment do not take it.
    \item Some assigned to control may obtain treatment anyway.
    \item Randomization applies to \(Z_i\), not necessarily to \(D_i\).
  \end{itemize}

  \begin{block}{Key shift}
    With non-compliance, random assignment becomes an instrument for treatment received.
  \end{block}
\end{frame}

\begin{frame}{Why Not Just Compare by Treatment Received?}
  Comparing \(Y_i\) by received treatment \(D_i\) can reintroduce selection:
  \[
    \E[Y_i\mid D_i=1]-\E[Y_i\mid D_i=0].
  \]

  \begin{itemize}
    \item People who comply may be richer, healthier, more motivated, or more informed.
    \item People who refuse treatment may differ in unobserved baseline outcomes.
    \item Treatment received is partly chosen, even if assignment was random.
  \end{itemize}

  \begin{block}{Example}
    If only users with new phones can receive a push notification, comparing delivered versus not delivered also compares richer versus poorer users.
  \end{block}
\end{frame}

\begin{frame}{Intention-to-Treat}
  Comparing by assignment is still causal:
  \[
    \ITT_Y
    =
    \E[Y_i\mid Z_i=1]-\E[Y_i\mid Z_i=0].
  \]

  \begin{itemize}
    \item This is the causal effect of being assigned or encouraged.
    \item It is often policy-relevant: what happens if we send offers?
    \item But it is diluted when many assigned people do not take treatment.
  \end{itemize}

  \begin{block}{Question}
    If we want the effect of actually receiving treatment, how do we undo the dilution?
  \end{block}
\end{frame}

\begin{frame}{Compliance Types}
  Define potential treatment received:
  \[
    D_i(1)=\text{treatment if assigned},
    \qquad
    D_i(0)=\text{treatment if not assigned}.
  \]

  \begin{center}
  \begin{tabular}{lccp{0.36\linewidth}}
    \toprule
    Type & \(D_i(1)\) & \(D_i(0)\) & Meaning \\
    \midrule
    Complier & 1 & 0 & follows assignment \\
    Never-taker & 0 & 0 & never takes treatment \\
    Always-taker & 1 & 1 & always takes treatment \\
    Defier & 0 & 1 & does the opposite of assignment \\
    \bottomrule
  \end{tabular}
  \end{center}
\end{frame}

\begin{frame}{Which Types Move the First Stage?}
  The first-stage effect of assignment on treatment is
  \[
    \E[D_i\mid Z_i=1]-\E[D_i\mid Z_i=0].
  \]

  Under random assignment, this equals
  \[
    \E[D_i(1)-D_i(0)].
  \]

  \begin{itemize}
    \item Compliers contribute \(1-0=1\).
    \item Never-takers contribute \(0-0=0\).
    \item Always-takers contribute \(1-1=0\).
    \item Defiers contribute \(0-1=-1\).
  \end{itemize}
\end{frame}

\begin{frame}{Monotonicity}
  \hypertarget{main:monotonicity}{}
  Monotonicity rules out defiers:
  \[
    D_i(1)\geq D_i(0)
    \quad \text{for every } i.
  \]

  \begin{itemize}
    \item Assignment can encourage treatment, or do nothing.
    \item It cannot make someone less likely to take treatment.
    \item This is often plausible for offers, eligibility, reminders, or draft risk.
  \end{itemize}

  \begin{block}{Why it matters}
    Without monotonicity, positive and negative first-stage responses can cancel, and the IV ratio is hard to interpret.
  \end{block}
\end{frame}

\section{LATE}

\begin{frame}{Local Average Treatment Effect}
  \hypertarget{main:late}{}
  Under independence, exclusion, relevance, and monotonicity:
  \[
    \frac{\E[Y_i\mid Z_i=1]-\E[Y_i\mid Z_i=0]}
         {\E[D_i\mid Z_i=1]-\E[D_i\mid Z_i=0]}
    =
    \E[Y_i(1)-Y_i(0)\mid D_i(1)>D_i(0)].
  \]

  \begin{block}{Name}
    This is the Local Average Treatment Effect, or LATE: the average treatment effect for compliers.
  \end{block}

  \derivbutton{late-proof}
\end{frame}

\begin{frame}{Why IV Identifies Compliers}
  The instrument only reveals treatment effects for people whose treatment status changes when \(Z_i\) changes.

  \begin{itemize}
    \item Never-takers are untreated regardless of assignment.
    \item Always-takers are treated regardless of assignment.
    \item Compliers switch treatment status when assignment changes.
  \end{itemize}

  \begin{block}{Core intuition}
    If the instrument does not change your treatment, your treatment effect cannot show up in the IV numerator.
  \end{block}
\end{frame}

\begin{frame}{LATE Is Not Automatically ATE}
  The ATE is
  \[
    \ATE=\E[Y_i(1)-Y_i(0)].
  \]
  LATE is
  \[
    \LATE=\E[Y_i(1)-Y_i(0)\mid \text{complier}].
  \]

  \begin{itemize}
    \item If treatment effects are constant, LATE equals ATE.
    \item If treatment effects differ by compliance type, LATE may be far from ATE.
    \item The IV estimate is internally valid for the complier group identified by that instrument.
  \end{itemize}

  \begin{block}{External validity question}
    Are the compliers the people policy makers care about?
  \end{block}
\end{frame}

\begin{frame}{Interpreting Compliers}
  Compliers are defined relative to the instrument.

  \begin{itemize}
    \item With a tuition subsidy, compliers are students induced to enroll by the subsidy.
    \item With quarter of birth, compliers are people whose schooling changes because school-entry and leaving rules bind.
    \item With a push notification assignment, compliers are users who receive the push only when assigned.
  \end{itemize}

  \begin{block}{Important}
    Different instruments for the same treatment can estimate different LATEs.
  \end{block}
\end{frame}

\begin{frame}{A Non-Compliance Numeric Example}
  Suppose an online course randomly offers free tutoring.

  \begin{center}
  \begin{tabular}{lcc}
    \toprule
    & Z=1 & Z=0 \\
    \midrule
    \(\E[D\mid Z]\) & 0.60 & 0.10 \\
    \(\E[Y\mid Z]\) & 82 & 77 \\
    \bottomrule
  \end{tabular}
  \end{center}

  \[
    \ITT_D = 0.60-0.10=0.50,
    \qquad
    \ITT_Y = 82-77=5.
  \]
  \[
    \widehat{\LATE}=\frac{\ITT_Y}{\ITT_D}=\frac{5}{0.50}=10.
  \]

  The offer raises scores by 5 points; tutoring raises scores by 10 points for students induced into tutoring by the offer.
\end{frame}

\begin{frame}{What To Report in an IV Paper}
  A credible IV table should not show only the second-stage coefficient.

  \begin{itemize}
    \item State the instrument and the treatment it shifts.
    \item Report first-stage coefficient and \(F\)-statistic.
    \item Report reduced form when useful.
    \item Explain independence and exclusion in institutional detail.
    \item State the likely complier population.
    \item Discuss why LATE is or is not the policy-relevant estimand.
  \end{itemize}
\end{frame}

\section{Simulation}

\begin{frame}{Python Example for This Lecture}
  The script for this lecture is:
  \[
    \texttt{lectures/code/lecture-5-iv-late.py}.
  \]

  It simulates four lessons:
  \begin{enumerate}
    \item OLS is biased when treatment is correlated with an omitted variable.
    \item IV recovers the treatment effect when the instrument is valid and strong.
    \item Weak first stages make IV estimates unstable.
    \item With non-compliance, Wald IV recovers LATE for compliers.
  \end{enumerate}

  \begin{block}{Why simulation helps}
    We know the true data-generating process, so students can see which estimator fails and why.
  \end{block}
\end{frame}

\begin{frame}{Simulation Design: Endogeneity and IV}
  The script generates
  \[
    U_i \sim N(0,1),
    \qquad
    Z_i \sim \text{Bernoulli}(1/2),
  \]
  \[
    D_i = 0.8Z_i + 0.7U_i + v_i,
  \]
  \[
    Y_i = 2D_i + 1.5U_i + e_i.
  \]

  \begin{itemize}
    \item True treatment effect is 2.
    \item OLS is biased because \(D_i\) contains \(U_i\).
    \item IV uses only the variation in \(D_i\) caused by \(Z_i\).
  \end{itemize}
\end{frame}

\begin{frame}{Simulation Design: Non-Compliance}
  For LATE, the script creates compliance types:
  \[
  (D_i(1),D_i(0))\in\{(1,0),(0,0),(1,1)\}.
  \]

  \begin{itemize}
    \item Compliers follow assignment.
    \item Never-takers never receive treatment.
    \item Always-takers receive treatment regardless of assignment.
    \item No defiers are generated, so monotonicity holds.
  \end{itemize}

  Treatment effects are heterogeneous, so the complier effect differs from the population ATE.
\end{frame}

\section{Review}

\begin{frame}{What Students Must Remember}
  \begin{enumerate}
    \item IV is for endogeneity from unobserved confounding, reverse causality, or measurement error.
    \item A valid instrument must be relevant, independent, and excluded from the outcome equation.
    \item Wald IV is reduced form divided by first stage.
    \item 2SLS operationalizes the same idea with controls and multiple instruments.
    \item Weak instruments are dangerous even when exclusion is conceptually plausible.
    \item With non-compliance and monotonicity, IV estimates LATE for compliers.
  \end{enumerate}
\end{frame}

\begin{frame}{Common Exam Mistakes}
  \begin{itemize}
    \item Saying ``\(Z\) is correlated with \(Y\), so it is a good instrument''.
    \item Forgetting that \(Z\) must move \(D\), not just predict \(Y\).
    \item Treating exclusion as a statistical test.
    \item Interpreting LATE as the ATE for everyone.
    \item Ignoring the first stage.
    \item Using treatment received as if it were randomized in a non-compliance setting.
  \end{itemize}
\end{frame}

\begin{frame}{Source and Further Reading}
  This lecture is based on ideas from:
  \begin{itemize}
    \item Matheus Facure Alves, \emph{Causal Inference for the Brave and True}, Chapter 08, ``Instrumental Variables''.
    \item Matheus Facure Alves, \emph{Causal Inference for the Brave and True}, Chapter 09, ``Non Compliance and LATE''.
    \item Angrist and Pischke, \emph{Mostly Harmless Econometrics}.
    \item Angrist and Pischke, \emph{Mastering 'Metrics}.
  \end{itemize}

  \vspace{0.3em}
  Handbook URLs:
  \begin{itemize}
    \item \url{https://matheusfacure.github.io/python-causality-handbook/08-Instrumental-Variables.html}
    \item \url{https://matheusfacure.github.io/python-causality-handbook/09-Non-Compliance-and-LATE.html}
  \end{itemize}
\end{frame}

\appendix

\begin{frame}{Appendix Map}
  \begin{itemize}
    \item A1: IV moment condition from the structural equation
    \item A2: Wald estimator derivation
    \item A3: 2SLS projection formula
    \item A4: IV with controls by residualizing
    \item A5: Why weak instruments inflate variance
    \item A6: LATE derivation under monotonicity
    \item A7: Compliance shares from first-stage moments
  \end{itemize}
\end{frame}

\begin{frame}{Appendix A1: IV Moment Condition}
  \hypertarget{app:iv-moment}{}
  \small
  Start from the structural equation
  \[
    Y_i = \alpha + \beta D_i + u_i.
  \]
  OLS fails when
  \[
    \E[D_i u_i]\neq 0.
  \]
  A valid instrument satisfies
  \[
    \E[Z_i u_i]=0,
    \qquad
    \E[Z_iD_i]\neq 0
  \]
  after demeaning or after including an intercept.

  Multiply the structural equation by \(Z_i\) and take expectations:
  \[
    \E[Z_iY_i]=\alpha\E[Z_i]+\beta\E[Z_iD_i]+\E[Z_iu_i].
  \]
  With de-meaned variables, \(\E[Z_i]=0\) and \(\E[Z_iu_i]=0\), so
  \[
    \beta=\frac{\E[Z_iY_i]}{\E[Z_iD_i]}.
  \]
  \backbutton{iv-conditions}
\end{frame}

\begin{frame}{Appendix A2: Wald Estimator}
  \hypertarget{app:wald-proof}{}
  For a binary instrument,
  \[
    \Cov(Y_i,Z_i)
    =
    \Prob(Z_i=1)\Prob(Z_i=0)
    \left(\E[Y_i\mid Z_i=1]-\E[Y_i\mid Z_i=0]\right).
  \]
  Similarly,
  \[
    \Cov(D_i,Z_i)
    =
    \Prob(Z_i=1)\Prob(Z_i=0)
    \left(\E[D_i\mid Z_i=1]-\E[D_i\mid Z_i=0]\right).
  \]
  Therefore the probability factor cancels:
  \[
    \frac{\Cov(Y_i,Z_i)}{\Cov(D_i,Z_i)}
    =
    \frac{\E[Y_i\mid Z_i=1]-\E[Y_i\mid Z_i=0]}
         {\E[D_i\mid Z_i=1]-\E[D_i\mid Z_i=0]}.
  \]
  \backbutton{wald}
\end{frame}

\begin{frame}{Appendix A3: 2SLS Projection Formula}
  \hypertarget{app:2sls-proof}{}
  \small
  Let \(X\) include exogenous controls and the intercept, \(D\) be the endogenous treatment, and \(Z\) include excluded instruments plus \(X\).

  The first-stage projection matrix is
  \[
    P_Z = Z(Z'Z)^{-1}Z'.
  \]
  The fitted treatment is
  \[
    \hat{D}=P_ZD.
  \]
  The 2SLS coefficient on \(D\) can be written as
  \[
    \hat{\beta}_{2SLS}
    =
    (D'P_ZM_XD)^{-1}D'P_ZM_XY,
  \]
  where \(M_X\) removes the included exogenous controls.

  In the just-identified no-control case, this collapses to
  \[
    \hat{\beta}_{2SLS}
    =
    \frac{Z'Y}{Z'D},
  \]
  after demeaning.
  \backbutton{2sls}
\end{frame}

\begin{frame}{Appendix A4: IV with Controls}
  \hypertarget{app:controls-proof}{}
  Suppose the model is
  \[
    Y_i = \alpha + \beta D_i + X_i'\gamma + u_i.
  \]
  Let \(\tilde{Y}_i\), \(\tilde{D}_i\), and \(\tilde{Z}_i\) be residuals from separate regressions of \(Y_i\), \(D_i\), and \(Z_i\) on \(X_i\).

  By Frisch-Waugh-Lovell, the controlled IV estimand is
  \[
    \beta_{IV\mid X}
    =
    \frac{\Cov(\tilde{Z}_i,\tilde{Y}_i)}
         {\Cov(\tilde{Z}_i,\tilde{D}_i)}.
  \]

  \begin{itemize}
    \item The numerator is the reduced form after removing \(X_i\).
    \item The denominator is the first stage after removing \(X_i\).
    \item The instrument must be valid conditional on \(X_i\).
  \end{itemize}
  \backbutton{wald}
\end{frame}

\begin{frame}{Appendix A5: Weak-IV Variance Intuition}
  \hypertarget{app:weak-proof}{}
  \footnotesize
  In the just-identified de-meaned case,
  \[
    \hat{\beta}_{IV}
    =
    \frac{\widehat{\Cov}(Z_i,Y_i)}
         {\widehat{\Cov}(Z_i,D_i)}.
  \]
  Let
  \[
    \pi = \frac{\Cov(Z_i,D_i)}{\Var(Z_i)}
  \]
  be the first-stage slope. When \(\pi\) is close to zero, the denominator is close to zero.

  A rough delta-method approximation gives
  \[
    \Var(\hat{\beta}_{IV}) \propto \frac{1}{\pi^2}.
  \]

  \begin{block}{Consequence}
    Halving the first-stage strength can roughly quadruple the variance, before considering finite-sample bias.
  \end{block}
  \backbutton{weak}
\end{frame}

\begin{frame}{Appendix A6: LATE Derivation}
  \hypertarget{app:late-proof}{}
  Let \(D_i(z)\) be potential treatment under assignment \(z\), and \(Y_i(d)\) potential outcome under treatment \(d\).

  Observed treatment:
  \[
    D_i = Z_iD_i(1)+(1-Z_i)D_i(0).
  \]
  By exclusion, observed outcome is
  \[
    Y_i = Y_i(0)+D_i\{Y_i(1)-Y_i(0)\}.
  \]
  By independence,
  \[
    \E[Y_i\mid Z_i=1]-\E[Y_i\mid Z_i=0]
    =
    \E\left[(D_i(1)-D_i(0))(Y_i(1)-Y_i(0))\right].
  \]
  With monotonicity, \(D_i(1)-D_i(0)\in\{0,1\}\), so the numerator equals
  \[
    \Prob(\text{complier})\E[Y_i(1)-Y_i(0)\mid \text{complier}].
  \]
  \backbutton{late}
\end{frame}

\begin{frame}{Appendix A6 Continued: Denominator}
  The first stage is
  \[
    \E[D_i\mid Z_i=1]-\E[D_i\mid Z_i=0]
    =
    \E[D_i(1)-D_i(0)].
  \]
  Under monotonicity,
  \[
    \E[D_i(1)-D_i(0)] = \Prob(\text{complier}).
  \]
  Therefore
  \[
    \frac{\E[Y_i\mid Z_i=1]-\E[Y_i\mid Z_i=0]}
         {\E[D_i\mid Z_i=1]-\E[D_i\mid Z_i=0]}
    =
    \E[Y_i(1)-Y_i(0)\mid \text{complier}].
  \]
  \backbutton{late}
\end{frame}

\begin{frame}{Appendix A7: Compliance Shares}
  \hypertarget{app:compliance-shares}{}
  Under random assignment:
  \[
    \E[D_i\mid Z_i=1]=\E[D_i(1)],
    \qquad
    \E[D_i\mid Z_i=0]=\E[D_i(0)].
  \]

  If monotonicity holds:
  \[
    \E[D_i\mid Z_i=1]-\E[D_i\mid Z_i=0]
    =
    \Prob(\text{complier}).
  \]

  If there are no always-takers, then
  \[
    \E[D_i\mid Z_i=0]=0,
    \qquad
    \E[D_i\mid Z_i=1]=\Prob(\text{complier}).
  \]

  If always-takers exist, \(\E[D_i\mid Z_i=0]\) estimates their share.
  \backbutton{noncompliance}
\end{frame}

\end{document}
